\section{Premise Audit and Scope}
\label{v480:premise-audit}

The following qualifications apply before reading the inherited core.
The original assertions and proofs are retained for audit, not silently
re-certified. This section records conditional replacements for specific
known issues; it does not claim a complete scientific validation of the
monolith or repair every downstream use of the older assertions.

\subsection{MM/HC Spectral Identification}
\label{v480:spectral-premises}

Commutation with an orthogonal projector preserves its image but does
not make that image a single eigenspace. Absence of nontrivial Jordan
chains also does not imply a single eigenvalue. For example,
$M=\operatorname{diag}(1,2,3)$ and $C=\operatorname{diag}(1,1,0)$
commute while $M$ is not scalar on $\operatorname{im}C$.
Even scalar action does not identify the full spectral projector:
$M=\operatorname{diag}(2,2,3)$ and $C=\operatorname{diag}(1,0,0)$
omit one direction of the eigenvalue-two eigenspace.
Thus the inherited labels
\texttt{lem:diagonal-traversal-scalar-action} and
\texttt{thm:hc-spectral-projector} must not be used without the
missing support premises.

\begin{proposition}[A sufficient conditional MM/HC identification]
\label{v480:conditional-mm}
Let $M$ be a bounded self-adjoint operator on a complex Hilbert space,
let $C$ be a nonzero orthogonal projector, and suppose for a real
$\lambda_0$ that
\[
 MC=CM=\lambda_0 C,
 \qquad \ker(M-\lambda_0 I)\cap\ker C=\{0\}.
\]
Then $\operatorname{im}C=\ker(M-\lambda_0 I)$, and $C$ is the
orthogonal eigenprojection for that eigenvalue. These conditions
are sufficient hypotheses, not a derived identification of a native MM.
\end{proposition}
\begin{proof}
The scalar relation gives the image-to-kernel inclusion. Conversely,
if $Mx=\lambda_0x$, commutation implies that $(I-C)x$ lies in both
$\ker(M-\lambda_0I)$ and $\ker C$. The second assumption forces
$(I-C)x=0$, so $x=Cx$. Orthogonality identifies the projection onto
the resulting closed eigenspace.
\end{proof}

One must still supply the realized operator, domain, scalar action and
complement exclusion. A representation label, an invariant carrier or
an audit token does not supply these hypotheses.

\subsection{OP1: Inverse Diagnostic and Fixed Exponent}
\label{v480:exponent-premises}

For $g>1$, $0<\beta<1/\sqrt3$ and $\gamma_1>0$, set
\[
 A(g,\beta)=\frac{\gamma_1}{4\pi|\log(\beta\sqrt3)|},
 \qquad n(g,\beta)=\frac{\log A(g,\beta)}{\log g}.
\]
The relation $g^{n(g,\beta)}=A(g,\beta)$ is then an identity.
Substituting $\beta(g)=\sqrt{(g-1)/(g+1)}$ leaves an identity on
$1<g<2$; it does not predict $g$. The inherited inverse-exponent
and self-consistency passages must be read with this restriction,
not as an independent arithmetic derivation of the frame factor.

\begin{proposition}[Unique root under an independently fixed ansatz]
\label{v480:fixed-exponent}
Assume the exponent five and prefactor $1/\sqrt3$ independently,
and put $a=\gamma_1/(4\pi)$ with $0<a<27/20$. The equation
\[
 f(g)=\tfrac12\log3+\tfrac12\log\frac{g-1}{g+1}
                   +\frac{a}{g^5}=0
\]
has exactly one root in $(1,2)$.
\end{proposition}
\begin{proof}
Here $f'(g)=1/(g^2-1)-5a/g^6$. For $t=g^2>1$,
\[
 t^3-\tfrac{27}{4}(t-1)=(t-\tfrac32)^2(t+3)\geq0.
\]
Consequently $g^6/(g^2-1)\geq27/4>5a$, so $f'>0$.
Also $f(1+)=-\infty$ and $f(2)=a/32>0$. Continuity and strict
monotonicity give existence and uniqueness.
\end{proof}

At the stated first-zero input $\gamma_1=14.134725141734694$,
the model root is $g\simeq1.863662249889$ and
$\beta\simeq0.549175395418$. This is a conditional scalar-model
calculation. The exponent, prefactor and prime-lattice interpretation
still need independent justification; the result does not close OP1/4.

\subsection{OP5: Detector Evidence and Full Balance}
\label{v480:op5-premises}

The full $\Theta$ sum, its regularization, endpoint conventions and
remainder control are separate obligations. An empty truncated
frequency sum does not make the full oscillatory sum vanish.
The inherited quantitative passages are therefore retained as audit
material, not certified by their relocation or compilation.
The specific conjecture for the full sum
(Proposition~\ref{conj:theta-sum-B}) has since been evaluated with the
corrected explicit formula in layer v4.8.0-r05
(\S\ref{sec:v480-r05-corrected-explicit-formula}): the left-hand side is
$-0.1097\pm0.0009$ against the required $0.893$, so that conjecture is
false. This decides the conjecture, not the general obligation above;
OP5 Step~(B) has no valid form in the declared reading.

The later C5 fit score is another interface. Applying a transfer or
contraction estimate to it requires the native evaluated boundary,
exterior and mixed residual blocks, their fixed norm conventions,
their guards and their calibration. Synthetic residual vectors or a
fitted contraction constant do not recover this missing dictionary.
A sufficient bound such as $3/4$ is not a measured universal physical
threshold, and a positive residual below a detector tolerance is not
an exact all-readout zero.

\subsection{Carried Corrections from the v4.8.0 Line}
\label{v480:carried-corrections}

The published v4.8.0 line and the layers recorded after it are appended
in the support route (\S\ref{v480:support-release-layers} and
\S\ref{v480:support-post-release-layers}). Two of their repairs were
made inside the full monolith and are not visible in the preserved
fragments, which keep the inherited wording. They apply here as follows.

\paragraph{Unity theorem.}
The preserved Theorem~\ref{thm:unity} lists the gauge structure
$U(1)\times SU(2)\times SU(3)$ as item~(vi) of its unitary-equivalence
list, and its proof structure cites Theorem~\ref{thm:SM} as~(vi).
Under layer r85 (\S\ref{sec:r85-citation-repair-unity-scope-separation})
the equivalence list ends at item~(v). A Lie group is not a Hilbert-space
representation of the Weyl algebra, so no Stone--von Neumann equivalence
is asserted for it; the gauge statement stands as a separate consequence,
and Theorem~\ref{thm:SM} is unchanged.

\paragraph{Stone--von Neumann citation.}
The proof structure of the unity theorem and one passage on the
irreducible representation of the Weyl algebra cite the Stone--von Neumann step as Theorem~3.2 of
\cite{brendecke2026rh}. In that source Theorem~3.2 is the functional
equation; the step is Step~3 of Lemma~13.2, resting on Proposition~13.1
(r85). Since layer v4.8.0-r04
(\S\ref{sec:v480-r04-riemann-decoupling}) the Riemann source is no longer
used as a proof source, the Stone--von Neumann justification of the unity
theorem is withdrawn, and the equivalence of items (i)--(v) is carried as
a conjecture.

\paragraph{Registers.}
Statements of the inherited core are to be read with the registers of
layer v4.8.0-r09 (\S\ref{sec:v480-r09-p2-p3-reclassification}) and
their extensions in r17--r20. In particular no prediction of the
inherited falsifiability corollary stands as confirmed. For the frame
layer: the frame factor used to convert constants between matter and
photon frame is the matter--tachyon cross term
(\S\ref{sec:v480-r18-three-frame-typing}), and in the complex-vector
reading no real Lorentz transformation relates the matter, tachyon and
photon sectors (\S\ref{sec:v480-r19-frame-cubes-holonomy},
\S\ref{sec:v480-r20-tachyon-complex-rapidity}). This agrees with the
present candidate's statement that no frame assignment is claimed.
